\documentclass{article} % This command is used to set the type of
% document you are working on such as an article, book, or presenation

\usepackage{geometry} % This package allows the editing of the page layout
\usepackage{amsmath}  % This package allows the use of a large range
% of mathematical formula, commands, and symbols
\usepackage{graphicx}  % This package allows the importing of images
\usepackage{amsfonts}

\newcommand{\question}[2][]{
  \begin{flushleft}
    \textbf{Question #1}: \textit{#2}

\end{flushleft}}
\newcommand{\sol}{\textbf{Solution}:} %Use if you want a boldface solution line
\newcommand{\maketitletwo}[2][]{
  \begin{center}
    \Large{\textbf{Assignment #1}

    Topology} % Name of course here
    \vspace{5pt}

    \normalsize{N. Ohayon Rozanes  % Your name here

    \today}        % Change to due date if preferred
    \vspace{15pt}

\end{center}}
\begin{document}
\maketitletwo[1]  % Optional argument is assignment number
%Keep a blank space between maketitletwo and \question[1]

\question[1]{}
\begin{enumerate}
  \item $$A \setminus (B \cap C) = (A \setminus B) \cup (A \setminus C)$$
    By definition we have that
    \begin{align*}
      A \setminus (B\cap C) &= \left\{x \vert x \in A \setminus (B
      \cap C) \right\} \\
      &= \left\{x \in A \text{ and } x \notin B \cap C \right\} \\
      &= \left\{x \in A \text{ and } x \notin B \text{ or }  x
      \notin C \right\}
    \end{align*}
    Similarily
    \begin{align*}
      (A \setminus B) \cup (A \setminus C) &= \left\{x \vert x \in
      (A\setminus B) \cup (A \setminus C)\right\} \\
      &= \left\{x | x \in A\setminus B \text{ or }  x \in A
      \setminus C \right\} \\
      &= \left\{x | x \in A\setminus B \text{ or }  x \in A
      \setminus C \right\} \\
      &= \left\{x | x \in A \text{ and } x \notin B \text{ or }  x \in A
      \text{ and } x \notin C \right\} \\
      &= \left\{x | x \in A \text{ and } x \notin B \text{ or }
      x \notin C \right\} \\
    \end{align*}
    As such,$\forall x \in A \setminus (B \cap C)$ it must be that
    $x \in A$ as well as $x \notin B$ or $x \notin C$, and therefore
    $x \in (A \setminus B) \cup (A \setminus C)$ as such we have that
    $A \setminus (B \cap C) \subset (A \setminus B) \cup (A \setminus
    C)$. Through a similar line of reasoning, $\forall x \in
    (A\setminus B) \cup (A \setminus C)$, we have that $x \in A$ as
    well as $x \notin B$ or $x \notin C$, so $(A \setminus B)
    \cup (A \setminus C)\subset A \setminus (B\cap C)$. Since by
    definition, two sets containing each other are equal, we have
    that $A \setminus (B \cap C) = (A \setminus B) \cup (A\setminus
    C)$ as desired.
  \item  $$ A \setminus (B \cup C) = (A \setminus B) \cap (A \setminus C)$$
    Using a similar process to the last problem, we have the
    following from definition
    \begin{align*}
      A \setminus (B \cup C) &= \{x | x \in A \text{ and } x \notin B
      \cup C \} \\
      &= \{x | x \in A \text{ and } x \notin B \text{ and } x \notin C \}
    \end{align*}
    And
    \begin{align*}
      (A \setminus B) \cap (A \setminus C) &= \{x | x \in A \setminus
      B \text{ and }  x \in A \setminus B \} \\
      &= \{x | x \in A \text{ and } x \notin B \text{ and } x \notin C\}
    \end{align*}
    Therefore $\forall x \in A \setminus (B \cup C)$, it follows that
    $x \in (A \setminus B) \cap (A \setminus C)$, so $A \setminus (B
    \cup C) \subset (A \setminus B) \cap (A \setminus C)$. In the
    same way $\forall x \in (A\setminus B) \cap (A \setminus C)$ it
    follows that $x \in A \setminus (B\cup C)$, therefore $(A
    \setminus B) \cap (A\setminus C) \subset A \setminus (B \cup C)$
    as such, we have that $A \setminus (B \cup C) = (A \setminus B)
    \cap (A \setminus C)$
\end{enumerate}

\question[2]{}
\begin{enumerate}
  \item Begin by considering the definition of the two sets
    $$
    \bigcap_{i \in I} A_i \cup \bigcap_{j\in J} B_i = \{x~|~x \in A_i
    \forall i \in I\} \cup \{x~|~x \in B_j \forall j \in J\}
    $$
    \[
      \bigcap_{i\in I, j\in J} (A_i \cup B_j) = \{x | \forall (i, j)
      \in I \times J, x \in A_i \cup B_j\}
    \]
    For an element $x$ to be included in the left hand side, it must be
    that its either in all $A_i$ or in all $B_j$, assume that $x \in
    A_i$ for all $i \in I$, then because $A_i  \subset A_i \cup B_j$,
    then  $x \in A_i \cup B_j$ $\forall i \in I$ or $j \in J$. Then
    it must be that $x$ is also in the right hand set. By symmetry,
    assume that $x\in B_j, \forall j \in J$, then it follows that $x
    \in A_i \cup B_j, \forall i \in I, j \in J$, therefore it follows that
    $$
    \bigcap_{i \in I} A_i \cup \bigcap_{j \in J} B_j \subset
    \bigcap_{i\in I, j\in J}
    (A_i \cup B_j)
    $$
    We now seek to show the reverse direction. Let $x \in
    \bigcap_{(i,j)} A_i \cup B_j$ , then either $\forall i
    \in I, x \in A_i$, in which case $x \in \bigcap A_i \subset
    \bigcap A_i \cup \bigcap B_j$ and we are done, or $\exists i_0
    \in I$ such that $x \notin A_{i_0}$. In which case, $x \in
    \bigcap A_i \cup B_j \implies x \in \bigcap B_j \implies x \in
    \bigcap B_j \cup \bigcap A_i$
    $$
    \bigcap_{i\in I, j\in J}
    (A_i \cup B_j)
    \subset
    \bigcap_{i \in I} A_i \cup \bigcap_{j \in J} B_j
    $$
    And therefore:
    $$
    \bigcap_{i\in I, j\in J}
    (A_i \cup B_j)
    =
    \bigcap_{i \in I} A_i \cup \bigcap_{j \in J} B_j
    $$
  \item
    \[
      \bigcup_{i \in I} A_i \cap \bigcup_{j \in J} B_j = \{x |
      \exists i \in I, x \in A_i \} \cap \{x | \exists j \in J, x \in B_j \}
    \]
    \[
      \bigcup_{i \in I,j \in J} A_i \cap B_j = \{ x| \exists i \in I
      \text{ and } j \in J, x \in A_i \cap B_j \}
    \]
    For any element $x$ to be in the left hand set, it must be that
    $\exists i_0 \in I$ such that $x \in A_{i_0}$ and also that $\exists j_0
    \in J$ such that $x \in B_{j_0}$, equivalently, $x \in A_{i_0}
    \cap B_{j_0}$. Then, $$x \in A_{i_0} \cap B_{j_0} \subset
    \bigcup_{i,j} A_i \cap B_j  \implies x \in \bigcup_{i,j} A_i \cap B_j$$
    And therefore
    \[\bigcup_{i \in I} A_i \cap \bigcup_{j \in J} B_j  \subset
      \bigcup_{i \in I,j \in J} A_i \cap B_j
    \]
    For the reverse direction, if an element $x$ is in the right hand
    set, there must be $i_0 \in I$ and $j_0 \in J$ such that $x \in
    A_{i_0} \cap B_{j_0}$. Since $$A_{i_0} \cap B_{j_0} \subset
    A_{i_0} \subset \bigcup_{i\in I} A_i$$ as well as $$A_{i_0} \cap B_{j_0}
    \subset B_{j_0} \subset \bigcup_{j \in J} B_j$$ it follows that
    \[
      x \in A_{i_0} \cap B_{j_0}  \subset \bigcup_{i \in I} A_i \cap
      \bigcup_{j \in J} B_j
    \]
    Since this holds for any element $x$ in the left hand set, we
    finally have that
    \[
      \bigcup_{i \in I,  j \in J} A_i \cap B_j \subset \bigcup_{i \in
      I} A_i \cap \bigcup_{j \in J} B_j
    \]
    \[
      \bigcup_{i \in I,  j \in J} A_i \cap B_j = \bigcup_{i \in
      I} A_i \cap \bigcup_{j \in J} B_j
    \]
    As desired.
\end{enumerate}

\question[3]{}
\begin{enumerate}
  \item Let $A \subset \mathbb{R}$ and $B \subset \mathbb{R}$, then
    $C = A \times B \subset \mathbb{R} \times \mathbb{R}$. Assume that $(a_1,
    b_1) \in C$ and $(a_2, b_2) \in C$. By definition:
    \[
      C = A\times B = \{(x,y)| x \in A \text{ and } y \in B\}
    \]
    That is, for any elements $a \in A$ and $b \in B, (a,b) \in C$
    We have from the assumptions that
    \[
      (a_1, b_1) \in C \text{ and } (a_2, b_2 ) \in C \implies a_1, a_2 \in
      A \text{ and } b_1, b_2 \in B
    \]
    It follows:
    \[
      a_1 \in A \text{ and } b_2 \in B \implies (a_1, b_2) \in C
    \]
    And
    \[
      a_2 \in A \text{ and } b_1 \in B \implies (a_2, b_1) \in C
    \]
  \item  The previous result gives us that if $C = A \times B$ for
    $A,B$ arbitrary subsets of $\mathbb{R}$, then it must be that
    $$(a_1, b_1) \in C \text{ and } (a_2, b_2) \in C \implies (a_1, b_2) \in C
    \text{ and } (a_2, b_1) \in C$$
    Taking the contrapositive of that, we have that
    if $\exists (a_1,
    b_1), (a_2, b_2) \in C$ and either $(a_1, b_2) \notin C$ or
    $(a_2, b_1) \notin C$ then there are no subsets of $\mathbb{R},
    A, B$ such that $C = A \times B$. It suffices to find $(a_1,
    b_1)$ and $(a_2, b_2)$ in $C$, such that $(a_1, b_2)$ is not in $C$.

    Consider the points $(1, 1)$ and $(-1, 1.4)$. It's clear that
    $$
    3^1 + 1^4 = 4 < 5 \implies (1,1) \in C
    $$
    And that
    $$
    3^{-1} + (1.4)^4 \approx 4.1749 < 5 \implies (-1, 1.4) \in C
    $$
    However,
    \[
      3^1 + (1.4)^4 = 6.8416 \ge 5 \implies (1, 1.4) \notin C
    \]
    Therefore, there does not exist subsets $A, B$ such that $C = A \times B$

\end{enumerate}

\question[4]{}

For $C$ to be an equivalence relation, it must be reflexive,
transative, and symmetric.
\begin{itemize}
  \item reflexivity: From the definition of a function, we have that
    $a = b \implies f(a) = f(b)$, and through the reflexivity of
    equality, $a = a$, so $f(a) =
    f(b)$, and therefore $aCa$
  \item transativity: If $f(a) = f(b)$, and $f(b) = f(c)$, then
    through transitivity of equality, we have that $f(a) = f(c)$. As
    such if $aCb$ and $bCc$, it must be that $aCc$
  \item symmetry: If $f(a) = f(b)$ then through symmetry of
    equivalence $f(b) = f(a)$, then we have that if $aCb \implies bCa$
\end{itemize}
Thus, $C$ is an equivlance relation.

\question[5]{}
\begin{enumerate}
  \item We demonstrate reflexivity, transitivity, and symmetry.
    \begin{itemize}
      \item reflexivity: For any $x \in \mathbb{Q}$, we have that $x
        - x = 0$, further, $6$ divides $0$ trivially, so $xCx$
      \item transitivity: For any $xCy$ and $yCz$, we have that $x -
        y  = 0 \mod 6$ and $y - z = 0 \mod 6$, this implies that $x -
        y + y - z = 0 \mod 6 \implies x - z = 0 \mod 6 \implies xCz$.
        Therefore if $xCy, y Cz \implies xCz$
      \item symmetry:
        \begin{align*}
          xCy &\implies x - y = 0 \mod 6\\
          &\implies -(x - y) = -(0)\mod 6 \\
          &\implies y - x = 0 \mod 6 \\
          &\implies yCx
        \end{align*}
        Thus $xCy \implies yCx$ as desired.
    \end{itemize}
  \item fix $x \in \mathbb{Q}$, then, for $k \in \mathbb{Z}$, then
    for any $y$ such that
    \begin{align*}
      xCy &\implies x - y = 0 \mod 6 \\
      &\implies x - y =  6k \\
      &\implies y = x - 6k
    \end{align*}
    Therefore, the equivalence class for $x$ is given by:
    \[
      [x] = \{x - 6k | \forall k \in \mathbb{Z}\}
    \]
\end{enumerate}
\end{document}
